\documentclass[10pt,a4paper]{amsart}
\usepackage[margin=19mm]{geometry}
\usepackage{amsmath,amssymb,mathtools}
\usepackage[scr=rsfs,cal=euler]{mathalfa}
\usepackage{xfrac}
\usepackage[unicode,hidelinks,bookmarks=false]{hyperref}
\newcommand{\ie}{i.e.\ }
\newcommand{\cf}{cf.\ }
\newcommand{\resp}{respectively\ }
\newcommand{\Subsection}{\subsection}
\providecommand{\nobreakdash}{\nobreak\hbox{-}}
\setlength{\emergencystretch}{2em}
\multlinegap0pt
\usepackage{amssymb}
\usepackage{mathtools}
\usepackage{tikz}
\usepackage{enumerate}
\usepackage{csquotes}
\newtheorem{proposition}{Proposition}
\title{Aggregation functions as lax morphisms of quantales}
\author{Alejandro Fructuoso-Bonet, Jes\'us Rodr\'{\i}guez-L\'opez}
\date{Source: arXiv:2604.26775v1 (CC BY 4.0)}
\begin{document}
\maketitle

\noindent\emph{Source and licence.} Aggregation functions as lax morphisms of quantales. Version arXiv:2604.26775v1 (CC BY 4.0), distributed by arXiv under the Creative Commons Attribution 4.0 International licence, http://creativecommons.org/licenses/by/4.0/, as recorded in the arXiv licence metadata for this exact version and shown on the abstract page https://arxiv.org/abs/2604.26775v1. Full licence notice: \texttt{licenses/arxiv-2604.26775-CC-BY-4.0.txt}.\par\smallskip

\noindent\textit{Editorial scope.} Counted unit: the article's proposition prop:Fdelta (source0.tex lines 979--985). The article's complete original proof of it is delivered with it (source0.tex lines 987--1057, one \texttt{proof} environment of 71 lines). No mathematics is written, repaired or re-derived: the statement and the proof are the article's own lines, in the article's own notation, and no retained line is substituted or re-flowed.  No further source-local context is needed: every cross-reference of every retained line resolves inside the two delivered spans.  Nothing was omitted from any retained span, so the package carries no omission record.  No retained line contains a citation, so the wrapper carries no bibliography and no bibliography entry.  No retained line contains a figure environment, an image-inclusion command or drawing code, so no image is delivered and the package declares no asset dependency.  Editorial formatting only, in addition: an AMS \texttt{amsart} preamble that restates the article's own declarations and macros used by the retained lines, namely the environments \texttt{proposition} on separate counters, together with the standard packages those lines need.  The retained environments print in this document as Proposition 1 in order of appearance, whereas the article prints the same items with the numbering of its own full document.  The complete native source of the article as distributed in the arXiv e-print for this exact version is bundled unchanged as \texttt{sources/199421/source0.tex}.
	\begin{proposition}\label{prop:Fdelta}
		Let $\ast$ be a continuous t-norm and $F:([0,1]^I,\leq,\ast)\to([0,1],\leq,\ast)$ be a map. Define  $F_\Delta:\Delta_+^I(\ast)\to \Delta_+(\ast)$ 
		as 
		$$F_\Delta((f_i)_i)(t)=F((f_i(t))_i)$$ for every $(f_i)_i\in \Delta_+^I(\ast)$ and every $t\in [0,+\infty].$
		
		Then $F_\Delta$ is a lax morphism if and only if $F$ is a left-continuous lax morphism.
	\end{proposition}

	\begin{proof}
		Suppose that $F$ is a left-continuous lax morphism. We first check that $F_\Delta$ is well-defined, that is, $F_\Delta((f_i)_i)\in\Delta_+(\ast)$ for any $(f_i)_i\in \Delta_+^I(\ast).$ Given $t,s\in [0,+\infty]$ such that $t\leq s$ then $(f_i(t))_i\leq (f_i(s))_i$ since $f_i\in \Delta_+(\ast)$ for every $i\in I$. Hence, $F_\Delta((f_i)_i)(t)=F((f_i(t))_i)\leq F((f_i(s))_i)=F_\Delta((f_i)_i)(s)$ since $F$ is isotone. Therefore $F_\Delta((f_i)_i)$ is isotone. 
		
		Furthermore, for every $t\in[0,+\infty]$ and $i\in I,$ since $f_i\in\Delta_+(\ast)$ we have that $f_i(t)=\bigvee_{0\leq s<t} f_i(s).$ Since $F$ is left-continuous then 
		\begin{align*}
			F_\Delta((f_i)_i)(t)&=F((f_i(t))_i)=F\left(\left(\bigvee_{0\leq s<t} f_i(s)\right)_i\right)=\bigvee_{0\leq s<t} F((f_i(s))_i)\\&=\bigvee_{0\leq s<t} F_\Delta((f_i)_i)(s).\end{align*}
		Consequently, $F_\Delta((f_i)_i)\in\Delta_+(\ast).$
		
		
		Next, we check that $F_\Delta$ is a lax morphism of quantales. 	Let $(g_i)_i,(h_i)_i\in \Delta_+^I(\ast)$ such that $(g_i)_i\leq (h_i)_i.$ Since $F$ is isotone then $F_\Delta((g_i)_i)(t)=F((g_i(t))_i)\leq F((h_i(t))_i)=F_\Delta((h_i)_i)(t)$ for all 
		$t\in [0,+\infty]$ . Thus $F_\Delta$ is isotone.
		
		Let $g=(g_i)_i,h=(h_i)_i\in \Delta_+^I(\ast).$ Let $t\in [0,+\infty]$ and $r,s\in [0,+\infty]$ such that $r+s\leq t.$ Since $F$ is a lax morphism then $F(g(r))\ast F(h(s))\leq F(g(r)\ast h(s))$. Consequently
		\begin{align*}			
			(F_\Delta(g)\circledast F_\Delta(h))(t)&=\bigvee_{r+s\leq t} F_\Delta(g)(r)\ast F_\Delta(h)(s)=\bigvee_{r+s\leq t} F((g_i(r))_i)\ast F((h_i(s))_i)\\&\leq \bigvee_{r+s\leq t} F((g(r)_i\ast h_i(s))_i)
			\leq F\left(\left(\bigvee_{r+s\leq t} g_i(r)\ast h_i(s)\right)_i\right)\\&=F(((g_i\circledast h_i)(t))_i)\\
			&=F_\Delta\left(g\circledast h\right)(t),
		\end{align*}
		
		and since $t$ is arbitrary
		$$ F_\Delta(g)\circledast F_\Delta(h)\leq F_\Delta(g\circledast h).$$
		
		Finally, 
		$$F_\Delta((f_{0,1})_i)(t)=F((f_{0,1}(t))_i)=\begin{cases}
			F((1)_i)&\text{ if }t>0,\\
			F((0)_i)&\text{ if }t=0,
		\end{cases}=\begin{cases}
			1&\text{ if }t>0,\\
			0&\text{ if }t=0,
		\end{cases}=f_{0,1},$$
		where we have used that $F((1)_i)=1$ since $F$ is a lax morphism and $F((0)_i)=F(\bigvee \varnothing)=\bigvee F(\varnothing)=0$ since $F$ is left-continuous.
		
		Therefore $F_\Delta$ is a lax morphism.
		
		
		Conversely, suppose that $F_\Delta$ is a lax morphism. Given $\alpha\in [0,1]$ define $f_\alpha\in\Delta_+(\ast)$ as
		$$f_\alpha(t)=\begin{cases}\alpha&\text{ if }t>0,\\
			0&\text{ if }t=0.\end{cases}$$
		Let $(u_i)_i,(v_i)_i\in [0,1]^I$ such that $(u_i)_i\leq (v_i)_i$. Then $(f_{u_i})_i\leq (f_{v_i})_i.$ Since $F_\Delta$ is isotone then $F_\Delta((f_{u_i})_i)\leq F_\Delta((f_{v_i})_i)$ so \begin{align*}
			F_\Delta((f_{u_i})_i)(1)&\leq F_\Delta((f_{v_i})_i)(1),\\
			F((f_{u_i}(1))_i)&\leq	F((f_{v_i}(1))_i),\\
			F((u_i)_i)&\leq F((v_i)_i).
		\end{align*}
		Hence $F$ is isotone.
		
		Moreover, if $(u_i)_i,(v_i)_i\in [0,1]^I$ are arbitrary then 
		\begin{align*}
			(F_\Delta((f_{u_i})_i)\circledast F_\Delta((f_{v_i})_i))(2)&\leq F_\Delta((f_{u_i})_i\circledast(f_{v_i})_i )(2),\\
			\bigvee_{r+s\leq 2} F_\Delta((f_{u_i})_i)(r)\ast  F_\Delta((f_{v_i})_i)(s)&\leq F\left(\bigvee_{r'+s'\leq 2} (f_{u_i}(r'))_i\ast (f_{v_i}(s'))_i\right).
		\end{align*}
		Since $(f_{u_i}(t))_i=(u_i)_i$ and $(f_{v_i}(t))_i=(v_i)_i$ for every $t>0$, the above supremum is attained for any pair $0<r+s\leq 2$ so 
		\begin{align*}
			F_\Delta((f_{u_i})_i)(1)\ast  F_\Delta((f_{v_i})_i)(1)&\leq F((f_{u_i}(1))_i\ast (f_{v_i}(1))_i),\\
			F((u_i)_i)\ast F((v_i)_i)&\leq F((u_i)_i\ast (v_i)_i).
		\end{align*}
		Furthermore, since $F_\Delta$ is a lax morphism then $F_\Delta((f_{0,1})_i)=f_{0,1}.$ Hence $F((1)_i)=F((f_{0,1}(1))_i)=F_\Delta((f_{0,1})_i)(1)=f_{0,1}(1)=1$. Consequently, $F$ is a lax morphism.
		
		At last, let us check that $F$ is left-continuous. Let $(t_i)_i\in [0,1]^I.$ If $t_i=0$ for all $i\in I$, since $F_\Delta(f_0)\in\Delta_+(\ast)$ then $0=F_\Delta((f_0)_i)(0)=F((f_0(0))_i)=F((0)_i).$ Suppose now that there exists $j\in I$ with $t_j\neq 0.$ For each $i\in I$ define $h_i\in \Delta_+(\ast)$ as
		$$h_i(t)=\begin{cases}
			t_i&\text{ if }t>1,\\
			t_i\cdot t&\text{ if }0\leq t\leq 1.
		\end{cases}$$
		Since $F_\Delta((h_i)_i)\in\Delta_+(\ast)$ then $F_\Delta((h_i)_i)(t)=\bigvee_{0\leq s<t}F_\Delta((h_i)_i)(s)$ for every $t\in [0,+\infty].$ In particular,
		\begin{align*}
			F_\Delta((h_i)_i)(1)&=\bigvee_{0\leq s<1}F_\Delta((h_i)_i)(r)\\
			F((h_i(1))_i)&=\bigvee_{0\leq r<1}F((h_i(r))_i)\\
			F((t_i)_i)&=\bigvee_{0\leq r<1}F((t_i\cdot r)_i)
		\end{align*}
		Since $\{(t_i\cdot r)_i:r\in [0,1)\}\subseteq \{(s_i)_i\in [0,1]^I: (0)_i\leq (s_i)_i<(t_i)\}$ and $F$ is isotone then 
		$$F((t_i)_i)=\bigvee_{0\leq r<1}F((t_i\cdot r)_i)\leq\bigvee_{(s_i)_i<(t_i)_i}F((s_i)_i)\leq F((t_i)_i) $$ so $F$ is left-continuous.
	\end{proof}

\end{document}
